\chapter{Count Your Knobs}

\epi{Not a verdict. A habit.}

\section{Four control panels}

We've now looked at four ways of explaining the sky, and counted the dials on each.

\fig{knob_counts}{The knob counts from this book, on a logarithmic scale: the law's one knob, \LCDM's six at the CMB, the fried chicken recipes from twelve to about 1066, and the string landscape somewhere far off the right edge of the page.}

\begin{table}[h]\centering\small
\begin{tabular}{@{}p{2.2cm}p{1.6cm}p{2.0cm}p{2.3cm}@{}}
\toprule
panel & knobs & covers & misses\\
\midrule
base \LCDM\ & 6 (+ nuisance) & CMB, expansion, light elements, large-scale structure & $\Lambda$'s value; the particle; galaxies need a second panel\\
\LCDM\ imitating the law & 12 to $\sim$1066, + 5 hand terms & the law's level and nearly its shape; clusters better than the law & the tie to dark energy; complete in one recipe\\
the landscape & $10^{500}$ clicks & everything, in principle & a map; established low-energy predictions\\
the law & 1 (+ shape, + scoped action) & disc galaxies, dwarfs, groups, lensing shape & clusters, CMB, several dwarfs, the Milky Way's vertical force\\
\bottomrule
\end{tabular}
\caption{The ledger. Nothing in it is a verdict on which theory is true.}
\end{table}

If you came for a winner, I'm sorry. The honest reading is that each panel is good at something and pays for it somewhere. \LCDM's six knobs are the most successful control panel in cosmology. Its galaxy panel can imitate my law closely, and with enough knobs it can even beat the law on clusters and lensing, but it can't make the law know about dark energy without typing it in. The landscape has so many clicks that it can't yet be checked. And my one knob describes disc galaxies with almost no freedom and fails, plainly and on the record, almost everywhere else.

\section{The six lessons}

Here's what the kitchen taught, compressed.

\begin{enumerate}
\item \textbf{The level is cheap.} Four knobs bring \LCDM\ halos to the law's acceleration scale within 3 per cent. A theory that reproduces ``about $10^{-10}\ms$'' hasn't done much.
\item \textbf{The shape is expensive.} Twelve knobs reach 0.0415\,dex, fourteen reach 0.0275, and stacking 28 ingredients plateaus near 0.022. The law's exact shape takes per-galaxy freedom or the law itself.
\item \textbf{The tie is out of reach.} No \LCDM\ route, of dozens tried, makes the acceleration scale follow the dark energy density without a hand term or a new coupling.
\item \textbf{An imitation proves freedom, not truth.} A tuned match, however close, is evidence that the model was flexible enough. It doesn't show the model is right, and it doesn't show the target is right either.
\item \textbf{Fewer knobs means more risk, not more truth.} The one-knob law fails in clusters and the CMB precisely because it has nowhere to hide. That's a feature of a falsifiable theory, and also a list of its failures.
\item \textbf{A knob that can't be turned is how a theory earns trust.} The law's prediction for how $a_0$ changes with redshift contains no galaxy parameter. It will come out right or wrong, and nobody can adjust it.
\end{enumerate}

\section{How to count the knobs in the next paper you read}

You don't need to agree with anything in this book to use its method. Next time you read a claim that a model ``explains'' something, ask six questions.

\begin{enumerate}
\item \textbf{What was frozen before the data were seen?} Targets, sample cuts, the statistic. If nothing was, the choice of what to report is itself a knob.
\item \textbf{How many numbers were set by the data being explained?} Count the per-object ones too. Six per galaxy over 171 galaxies is 1026.
\item \textbf{Are any knobs at their bounds?} A model pushing half its population to the edge of a prior is asking for different physics.
\item \textbf{How does it do on data it never saw?} A holdout grade counts for more than a tuning grade.
\item \textbf{Were any dependences written in by hand?} If so, the behaviour they produce is an input, not an output.
\item \textbf{What does it predict that it can't adjust?} If the answer is ``nothing,'' then whatever the model's merits, it hasn't yet stuck its neck out.
\end{enumerate}

Apply those to my law and you get: one fitted knob, a shape choice, an action adopted in a scope, a frozen scorecard with six failures out of eleven, and one unturnable prediction. Apply them to the fried chicken recipes and you get: 12 to 1066 knobs, several at their bounds, decent holdout grades, five hand terms for the tie, and no unturnable prediction, because the next sample can always be tuned to as well. Apply them to the landscape and you get: zero knobs on the panel, $10^{500}$ in the warehouse, and a selection problem bigger than all of physics' measurements put together.

\section{What happens next}

The calendar is short. On 6 November 2026, the day after this edition is dated, an ALMA observation of one high-redshift galaxy goes public, and the programme scores it against predictions frozen in advance. It tests one galaxy's mass budget, not the tie, and one disc excludes no model. Around December, Gaia's fourth data release gives the wide-binary test a much larger sample, with the law's prediction sealed at 1.149.

The real test of the knob that can't be turned isn't on any calendar yet. It needs a few galaxies at redshift two or more, measured cleanly in the deep regime, with stars and gas weighed independently of their rotation and a shared mass error under about 0.08\,dex. That's hard. It's within reach of the telescopes we have, and it's the kind of measurement that ends arguments.

When it's made, one of two things will happen. Either $a_0$ will have climbed with the expansion rate, like every halo scale in this book, and the law will be finished. Or it will have stayed put, following the dark energy history that the expansion surveys measure, and the law will have done the one thing no knob could do for it.

Either way, nobody will have turned a knob. That's the point.

\bigskip
\begin{center}
\emph{Count your knobs.}
\end{center}
